\chapter{Set up Git and GitHub for submission}
\label{ch:git-submission}

The \href{https://git-scm.com/install/}{\textbf{official Git download page}} lists current installers and package-manager commands.

\subsection{Install Git}

Use the command for your operating system, then reopen the terminal:

\begin{lstlisting}[style=utpterminal]
# Windows PowerShell
winget install --id Git.Git -e --source winget

# macOS (Option A): install using Homebrew
brew install git

# macOS (Option B): install Apple's command-line tools
xcode-select --install

# Linux
sudo apt update
sudo apt install git
\end{lstlisting}

On Windows, use \textbf{Git Bash} for the SSH commands below. Confirm the installation on any system:

\begin{lstlisting}[style=utpterminal]
git --version
\end{lstlisting}

\subsection{Configure your commit identity}

Set the name and email that should appear on future commits. Use an email
connected to your GitHub account, or your GitHub-provided private email.

\begin{lstlisting}[style=utpterminal]
git config --global user.name "Your Name"
git config --global user.email "you@example.com"
git config --global init.defaultBranch main

git config --global --list
\end{lstlisting}

\subsection{Video walkthroughs for SSH setup}

Use the walkthrough for your operating system if you prefer to follow the SSH
setup on screen:

\begin{tcolorbox}[
  enhanced,
  breakable,
  colback=UTPPaleOrange,
  colframe=UTPOrange,
  colbacktitle=UTPOrange,
  boxrule=1pt,
  arc=2mm,
  title={Choose your operating system},
  fonttitle=\bfseries\color{white},
  left=3mm,right=3mm,top=2mm,bottom=2mm
]
\begin{description}[leftmargin=1.15in,style=nextline]
  \item[Linux]
    \href{https://youtu.be/AOM8CEyu3Ug?t=121}{Set up GitHub SSH on Linux}
  \item[Windows]
    \href{https://youtu.be/YT5YimhyX18?t=5}{Set up GitHub SSH on Windows}
  \item[macOS]
    \href{https://youtu.be/ACvE5Uy7ZB0?t=437}{Set up GitHub SSH on macOS}
\end{description}
\end{tcolorbox}

\pagebreak
\subsection{Create an SSH key for GitHub}

SSH lets Git authenticate without asking for a GitHub password on every push.
First check whether this computer already has a key:

\begin{lstlisting}[style=utpterminal]
ls -al ~/.ssh
\end{lstlisting}

If \texttt{id\_ed25519} and \texttt{id\_ed25519.pub} do not exist, create them.
Enter the email connected to GitHub and choose a passphrase when prompted:

\begin{lstlisting}[style=utpterminal]
ssh-keygen -t ed25519 -C "you@example.com"
\end{lstlisting}

The command above will give:

\begin{enumerate}[leftmargin=1.7em]
  \item Generating public/private ed25519 key pair.
  \item Enter file in which to save the key (\textasciitilde/.ssh/id\_ed25519): \textit{(Press Enter to accept the default)}
  \item Enter passphrase (empty for no passphrase): \textit{(Press Enter for no passphrase)}
  \item Enter same passphrase again: \textit{(Press Enter to confirm)}
\end{enumerate}

\begin{lstlisting}[style=utpterminal]
Your identification has been saved in C:\Users\USERNAME\.ssh\id_ed25519
Your public key has been saved in C:\Users\USERNAME\.ssh\id_ed25519.pub
The key's randomart fingerprint is:
+--[ED25519 256]--+
|   ..oo+.        |
|   o.+.= .       |
|  . o o =        |
|   . E o         |
|    o . S        |
|     o o         |
|      o          |
+-----------------+
\end{lstlisting}

After that you can load the \emph{private} key into the SSH agent:

For Windows PowerShell:

\begin{lstlisting}[style=utpterminal]
Set-Service -Name ssh-agent -StartupType Automatic
Start-Service ssh-agent
ssh-add $env:USERPROFILE\.ssh\id_ed25519
git config --global core.sshCommand "C:/Windows/System32/OpenSSH/ssh.exe"
\end{lstlisting}

For macOS:
\begin{lstlisting}[style=utpterminal]
eval "$(ssh-agent -s)"
ssh-add --apple-use-keychain ~/.ssh/id_ed25519
\end{lstlisting}

For Linux:
\begin{lstlisting}[style=utpterminal]
eval "$(ssh-agent -s)"
ssh-add ~/.ssh/id_ed25519
\end{lstlisting}

Copy the \emph{public} key. Never upload or share the file without
\texttt{.pub}.

\begin{lstlisting}[style=utpterminal]
# Windows PowerShell
Get-Content $env:USERPROFILE\.ssh\id_ed25519.pub | Set-Clipboard

# macOS
pbcopy < ~/.ssh/id_ed25519.pub

# Linux: display it, then copy the complete line
cat ~/.ssh/id_ed25519.pub
\end{lstlisting}

\noindent To authenticate Git operations securely without entering your password, copy your public SSH key and register it in your GitHub account:

\begin{enumerate}[leftmargin=1.7em]
  \item Copy your public key directly to your clipboard. If you generated an ed25519 key, run the command below in PowerShell or terminal. Do not share or copy the private key (\texttt{id\_ed25519}).
  \item In the upper-right corner of GitHub, click your profile icon and select \textbf{Settings}.
  \item In the left sidebar under \textbf{Access}, select \textbf{SSH and GPG keys}.
  \item Click \textbf{New SSH key} (or \textbf{Add SSH key}).
  \item In the \textbf{Title} field, enter a descriptive label such as \texttt{personal laptop}.
  \item Ensure \textbf{Key type} is set to \textbf{Authentication Key}.
  \item Paste the public key string into the \textbf{Key} field.
  \item Click \textbf{Add SSH key}, then confirm your GitHub account password or two-factor prompt if requested.
\end{enumerate}

\noindent Verify that your machine successfully connects and authenticates with GitHub:

\begin{lstlisting}[style=utpterminal]
ssh -T git@github.com
\end{lstlisting}


\subsection{Prepare the course repository}

Do not upload the virtual environment, passwords, tokens, or API keys. In the project root, create a file named \texttt{.gitignore}.

Create a new repository on GitHub. Use a name such as \texttt{UTP-"COURSE-NAME"/}, select the visibility required by the instructor, and leave \textbf{README}, \textbf{.gitignore}, and \textbf{license} unchecked because the local project already exists.

From the project root, create the first commit and connect it to the empty
GitHub repository. Replace \texttt{USERNAME} and \texttt{REPOSITORY}:

\begin{lstlisting}[style=utpterminal]
git init -b main
git add .
git status
git commit -m "adding course content"

git remote add origin git@github.com:USERNAME/REPOSITORY.git
git remote -v
git push -u origin main
\end{lstlisting}

\begin{pitfall}
Always read \texttt{git status} before committing. If a dataset, password,
token, \texttt{.env} file, or virtual environment appears, stop and update
\texttt{.gitignore} before running \texttt{git commit}. Deleting a secret in a
later commit does not remove it from earlier repository history.
\end{pitfall}

\subsection{Submit later changes}

Use short commits as each lab progresses:

\begin{lstlisting}[style=utpterminal]
git status
git add .
git commit -m "adding lab..."
git push
\end{lstlisting}

Submit the GitHub repository address requested by the instructor and confirm that the latest commit and required notebooks are visible on GitHub.